\documentclass[a4paper,12pt]{article}
\usepackage{cmap}
\usepackage[T2A]{fontenc}
\usepackage[utf8]{inputenc}
\usepackage[english,russian]{babel}
\usepackage{amsmath,amsfonts,amssymb,amsthm,mathtools,mathtext}
\usepackage[left=1cm,right=1cm,top=2cm,bottom=2cm]{geometry}
\usepackage{graphicx}
\usepackage[colorlinks,linkcolor=blue]{hyperref}
\usepackage{physics}
\usepackage{esvect}
\usepackage{icomma}
\usepackage{indentfirst}
\usepackage[labelsep=period,justification=centering]{caption}
\usepackage{array}
\usepackage{multirow}
\usepackage{float}
\usepackage{longtable}
\usepackage{tabularray}
\usepackage{booktabs}
\usepackage{tabularx}

\title{Лабораторная работа 1.2.1. Определение скорости пули при помощи маятника}
\author{Aleksandr Chirkin}
\date{October 2026}

\begin{document}

\maketitle

\section{Аннотация}
В работе определяется скорость полёта пули при помощи баллистических маятников,
производится оценка погрешности измерений.

\section{Теоретические сведения}
\subsection{Баллистический маятник, совершающий поступательное движение}

\begin{figure}[h]
    \centering
    \includegraphics[width=0.75\textwidth]{setup1.pdf}
    \caption{Баллистический маятник, совершающий поступательное движение}
    \label{fig:setup1}
\end{figure}

В первом опыте использовалась установка, представленная на рисунке
\ref{fig:setup1}. Пусть масса маятника $M$, масса пули $m$, $m \ll M$. Тогда
скорость пули выражается по формуле 
\[
    u \approx \frac{M}{m}V
\]

Используя закон сохранения механической энергии, для малых углов отклонения
получим формулу скорости пули:
\begin{equation}
    \label{eq:bullet_speed1}
    u = \frac{M}{m}\sqrt{\frac{g}{L}}\Delta x
\end{equation}

\subsection{Баллистический маятник, совершающий вращательное движение}
\begin{figure}[h]
    \centering
    \includegraphics[width=0.75\textwidth]{setup2.pdf}
    \caption{Баллистический маятник, совершающий вращательное движение}
    \label{fig:setup2}
\end{figure}

Во втором опыте использовалась установка, представленная на рисунке
\ref{fig:setup2}. $r$ --- расстояние от линии полёта пули до оси вращения
маятника, $I$ --- момент инерции маятника, $\Omega$ --- его угловая скорость
сразу после удара. $M$ --- масса грузов, $m$ --- масса пули.
По законам сохранения энергии и момента импульса:
$$mur = I\Omega, \quad k\frac{\varphi^2}{2} = I\frac{\Omega^2}{2}$$
Так как угол $\varphi$ мал:
$$\varphi \approx \frac{x}{2d}$$
Период колебания маятника с грузами и без грузов:
$$T_1 = 2\pi\sqrt{\frac{I}{k}}, \quad T_2 = 2\pi\sqrt{\frac{I-2MR^2}{k}}$$

Окончательно получим:
\begin{equation}
    \label{eq:bullet_speed2}
    u = \frac{4\pi MR^2T_1x}{(T_1^2-T_2^2)2mrd}
\end{equation}

Кроме того, как и в любом другом эксперименте, нам необходимо учесть погрешность измерений:

Для поступательного маятника формула относительной погрешности выглядит следующим образом:
\begin{equation}
\varepsilon_u = 100\% \cdot \sqrt{ 
    \begin{aligned}[t]
        & \left(\frac{\Delta M}{M}\right)^2 + \left(\frac{\Delta m}{m}\right)^2 + \left(\frac{\Delta L}{2L}\right)^2 + {} \\
        & + \left(\frac{\Delta(\Delta x)}{\Delta x}\right)^2
    \end{aligned}
}
\end{equation}

Для вращательного маятника:
\begin{equation}
\begin{aligned}
\varepsilon_u = 100\% \cdot \sqrt{ 
    \begin{aligned}[t]
        & \left(\frac{\Delta M}{M}\right)^2 + \left(2\frac{\Delta R}{R}\right)^2 + \left(\frac{\Delta x}{x}\right)^2 + \left(\frac{\Delta m}{m}\right)^2 + {} \\
        & + \left(\frac{\Delta r}{r}\right)^2 + \left(\frac{\Delta d}{d}\right)^2 + \left(\frac{T_1^2 + T_2^2}{T_1(T_1^2 - T_2^2)} \Delta T_1\right)^2 + \left(\frac{2T_2}{T_1^2 - T_2^2} \Delta T_2\right)^2
    \end{aligned}
}
\end{aligned}
\end{equation}

\section{Измеренные данные}

\begin{table}[htbp]
    \centering
    \caption{Массы пуль}
    \label{tab:bullet_weights}
    \begin{tabular}{cc}
        \toprule
        Номер пули & Масса пули, г \\
        \midrule
        1 & 0.506 \\
        2 & 0.511 \\
        3 & 0.505 \\
        4 & 0.517 \\
        5 & 0.505 \\
        6 & 0.508 \\
        7 & 0.512 \\
        8 & 0.502 \\
        9 & 0.509 \\
        10 & 0.511 \\
        \bottomrule
    \end{tabular}
\end{table}

\subsection{Инерциальный баллистический маятник}
Для инерциального маятника $M = 2925 \pm 5 \text{ г}$, $L = 208.2 \pm 0.2 \text{ см}$.

\begin{table}[H]
    \centering
    \caption{Параметры пуль}
    \label{tab:bullets_data}
    \begin{tabular}{ccc}
        \toprule
        № пули & Амплитуда $\Delta x$, мм & Масса $m$, г \\
        \midrule
        1 & 12.0 & 0.506 \\
        2 & 12.0 & 0.511 \\
        3 & 11.5 & 0.505 \\
        4 & 12.0 & 0.517 \\
        5 & 11.4 & 0.505 \\
        \bottomrule
    \end{tabular}
\end{table}

Погрешности измеренных величин составляют: $\Delta(\Delta x) = 0.1 \text{ мм}$, $\Delta(m) = 0.001 \text{ г}$.
После подстановки в формулы (1) и (3) получаем результаты, представленные в таблице \ref{tab:bullet_speed_errors}.

\begin{table}[H]
    \centering
    \caption{Результаты расчета скорости и погрешностей}
    \label{tab:bullet_speed_errors}
    \begin{tabularx}{\textwidth}{c>{\centering\arraybackslash}X>{\centering\arraybackslash}X>{\centering\arraybackslash}X}
        \toprule
        № пули & Скорость $u$, м/с & Относительная погрешность $\varepsilon_u$, \% & Абсолютная погрешность $\Delta u$, м/с \\
        \midrule
        1 & 150.497 & 2.598 & 3.910 \\
        2 & 149.025 & 2.583 & 3.849 \\
        3 & 144.512 & 2.648 & 3.827 \\
        4 & 147.295 & 2.566 & 3.779 \\
        5 & 143.256 & 2.658 & 3.808 \\
        \bottomrule
    \end{tabularx}
\end{table}

По итогам можно установить, что скорость пули в опытах 1 - 5 была равна $u = 146.917 \pm 3.835 \text{ м/с}$.
\subsection{Крутильный баллистический маятник}

Для крутильного маятника $M = 1468,7 \pm 0.2 \text{ г}$, $R = 34 \pm 0.1\text{ см} $, $r = 19 \pm 0.1\text{ см} $, $d = 135 \pm 0.1\text{ см} $, $T_1 = 6,56 \pm 0.1 \text{ с}$, $T_2 = 4,85  \pm 0.1 \text{ с}$

\begin{table}[htbp]
    \centering
    \caption{Параметры пуль (продолжение, 6--10 опыты)}
    \label{tab:bullets_data_2}
    \begin{tabular}{ccc}
        \toprule
        № пули & Амплитуда $\Delta x$, см & Масса $m$, г \\
        \midrule
        6 & 3,9  & 0.508 \\
        7 & 4,0  & 0.512 \\
        8 & 4,1  & 0.502 \\
        9 & 4,0  & 0.509 \\
        10 & 4,0 & 0.511 \\
        \bottomrule
    \end{tabular}
Погрешности измеренных величин составляют: $\Delta(\Delta x) = 1 \text{ мм}$, $\Delta(m) = 0.001 \text{ г}$.
\end{table}

После подстановки в формулы (2) и (4) получаем результаты, представленные в таблице 

\begin{table}[htbp]
    \centering
    \caption{Результаты расчета скорости и погрешностей (продолжение, 6--10 опыты)}
    \label{tab:bullet_speed_errors_2}
    \begin{tabularx}{\textwidth}{c>{\centering\arraybackslash}X>{\centering\arraybackslash}X>{\centering\arraybackslash}X}
        \toprule
        № пули & Скорость $u$, м/с & Относительная погрешность $\varepsilon_u$, \% & Абсолютная погрешность $\Delta u$, м/с \\
        \midrule
        6 & 107.351 & 7.928 & 8.510 \\
        7 & 109.243 & 7.904 & 8.634 \\
        8 & 114.205 & 7.894 & 9.015 \\
        9 & 109.887 & 7.907 & 8.689 \\
        10 & 109.457 & 7.905 & 8.653 \\
        \bottomrule
    \end{tabularx}
\end{table}

По итогам можно установить, что скорость пули в опытах 6-10 была равна $u = 110,05 \pm 8,7 \text{ м/c}$

\section{Выводы}
В работе удалось произвести измерение скорости пули с точностью порядка 2,5\% на
баллисти\-ческом маятнике, совершающем поступательное движение, и с точностью
порядка 8\% на крутильном баллистическом маятнике. Наибольший вклад
в погрешность составляют погреш\-ности измерения расстояний, которые являются
высокими из-за особенностей проведения опыта, и их можно снизить. Разброс
значений не даёт возможности определить зависимость скорости пули от её массы,
для этого требуется проводить существенно более точные измерения.
\end{document}